\section{Original auxiliary labels in physical-copy coordinates}
\label{sec:regional_auxiliary_labels}

\begin{theorem}[The two original auxiliary projections]
\label{thm:global_replica_auxiliary_labels}
\lean{TensorPower.labelProj_globalReplicaCopiesEquiv_auxiliary}
Let $\mathcal H$ be the original finite physical tensor product, and let
$C,R$ be finite auxiliary coordinate spaces. For a copy number $k$, let
$G_k$ be the fixed regrouping of
Definition~\ref{def:regional_five_factor_coordinates}, with codomain
$\mathcal H^{\otimes k}\otimes C^{\otimes k}\otimes R^{\otimes k}$.
For any original labels $\ell_C,\ell_R$, let $P_C,P_R$ denote their
whole-copy auxiliary projections. Then
\begin{align}
 G_k^*(I\otimes P_C\otimes I)G_k&=P_{\{C\},\ell_C},\\
 G_k^*(I\otimes I\otimes P_R)G_k&=P_{\{R\},\ell_R}.
\end{align}
On the right, $C$ and $R$ denote the two distinguished original auxiliary
sites. The identities hold for zero copies and empty coordinate spaces;
no state or regional cut is needed. They give the auxiliary-coordinate
identification used in
\cite[Section~8, lines~421--443]{OpenAI2026AreaLaw}.
\end{theorem}
\begin{proof}
\uses{def:regional_five_factor_coordinates}
The regrouping intertwines the permutation of the $C$ copies with the
same permutation of the original $C$ site, and similarly for $R$.
Taking the central idempotent of each label gives the two identities.
The other coordinate factors carry the identity.
\end{proof}
